\chapter{Reading Query Plans}
\label{app:plans}

A reference for the physical-plan operators the book refers to. \Cref{ch:catalyst}
explains the compilation pipeline that produces these plans and how to read one
quickly; this appendix is the catalogue you come back to. Everything here is as
of the reference environment (\Cref{app:refenv}).

\section{The six \texttt{EXPLAIN} modes}

\begin{table}[H]
\centering\footnotesize
\renewcommand{\arraystretch}{1.15}
\begin{tabular}{@{}p{3cm}p{10.8cm}@{}}
\toprule
\textbf{Form} & \textbf{Output} \\
\midrule
\texttt{EXPLAIN} & physical plan only \\
\texttt{EXPLAIN EXTENDED} & parsed, analyzed, optimized logical, and physical plans \\
\texttt{EXPLAIN FORMATTED} & physical plan as a tree + a numbered list expanding every node \\
\texttt{EXPLAIN COST} & optimized logical plan with the planner's row-count / size estimates \\
\texttt{EXPLAIN CODEGEN} & physical plan + the generated Java for each whole-stage stage \\
\texttt{df.explain(mode=)} & same five, as \texttt{"simple"}, \texttt{"extended"}, \texttt{"formatted"}, \texttt{"cost"}, \texttt{"codegen"} \\
\bottomrule
\end{tabular}
\end{table}

\section{Reading a physical plan}

Read \textbf{bottom-up}: the leaves are scans, the root is the last operator to
run. The tree connectors are \texttt{+-} for the child in line and \texttt{:-}
plus a \texttt{:} spine for the first child of a binary operator (a join). A
\texttt{*} prefix, written \texttt{*(n)}, marks an operator fused into
whole-stage-generated Java stage \texttt{n} (\Cref{sec:c-codegen}).

Here is that process actually applied, to the specimen's own
\texttt{orders}~$\bowtie$~\texttt{dim\_product} plan from
\Cref{lst:explain-anatomy}:

\begin{lstlisting}[language=text,numbers=none,breaklines=false,basicstyle=\footnotesize\ttfamily]
*(3) Sort (13)
+- *(3) HashAggregate (12)
   +- *(2) HashAggregate (10)
      +- *(2) SortMergeJoin Inner (9)
         :- *(1) Sort (5)
         :  +- Exchange hashpartitioning(product_id, 200) (4)
         :     +- *(1) Filter (3)
         :        +- *(1) ColumnarToRow (2)
         :           +- BatchScan orders (1)
         +- *(1) Sort (8)
            +- Exchange hashpartitioning(product_id, 200) (7)
               +- BatchScan dim_product (6)
\end{lstlisting}

Reading it bottom-up, one operator at a time: node~(1) is the first thing
that runs --- \texttt{BatchScan orders} reads the raw files. Node~(2)
converts the columnar batches it produces into row form. Node~(3) applies
the \texttt{order\_ts} range filter, discarding rows outside the query's
week. That is the entire left-hand side up to node~(4), an
\texttt{Exchange}: the filtered rows are shuffled by
\texttt{hash(product\_id)} into 200 partitions, so that every row with a
given \texttt{product\_id} ends up in the same partition as every other row
sharing it. Node~(5) locally sorts each of those partitions by
\texttt{product\_id}. The right-hand branch, nodes~(6)--(8), does the exact
same three things --- scan, shuffle, sort --- to \texttt{dim\_product},
independently and in parallel with the left branch. Only once both branches
have reached a sorted, co-partitioned state does node~(9),
\texttt{SortMergeJoin}, run: it walks both sorted streams together,
matching rows with equal \texttt{product\_id}. From there, node~(10) runs a
partial aggregation per partition (no new exchange needed, since the join
already grouped matching keys together), node~(12) finishes that
aggregation, and node~(13) sorts the small, already-aggregated result for
final output. Every step in this reading is exactly "what does this node
need before it can run, and has that already happened lower in the tree" ---
which is the whole discipline bottom-up reading is.

\section{Scan operators}
\label{sec:c-scan}

Where data enters the plan.

\begin{table}[H]
\centering\footnotesize
\renewcommand{\arraystretch}{1.15}
\begin{tabular}{@{}p{3.6cm}p{10.2cm}@{}}
\toprule
\textbf{Operator} & \textbf{Source} \\
\midrule
\texttt{BatchScan} & a DataSource~V2 table --- Iceberg on the reference stack \\
\texttt{FileScan parquet} / \texttt{orc} / \texttt{csv} & a file-based table through the Hive catalog (DataSource~V1) \\
\texttt{InMemoryTableScan} & a cached DataFrame or \texttt{CACHE TABLE} \\
\texttt{DataSourceV2ScanRelation} & the logical-plan form of a V2 scan \\
\bottomrule
\end{tabular}
\end{table}

\begin{lstlisting}[language=text,numbers=none,caption={A scan node's fields},captionpos=b]
BatchScan orders
  PartitionFilters: [order_ts >= 2026-03-01, order_ts < 2026-03-08]
  PushedFilters:    [IsNotNull(product_id)]
  DataFilters:      [IsNotNull(product_id)]
  ReadSchema:       struct<product_id,promised_days,delivery_days,order_ts>
  Batched:          true
\end{lstlisting}

\begin{description}
  \item[\texttt{PartitionFilters}] filters resolved against partition values, so
        whole partitions (directories on object storage) are skipped without
        being opened. Empty on a partitioned table means the filter did not match
        the partition transform (\Cref{ch:pruning}).
  \item[\texttt{PushedFilters}] filters handed to the file reader, which skips
        row groups using column min/max statistics in the footer and any bloom
        filters, without decompressing them.
  \item[\texttt{DataFilters}] every filter on the scan, including the residue
        evaluated row-by-row after reading. \texttt{DataFilters} minus
        \texttt{PushedFilters} is what could not be pushed.
  \item[\texttt{ReadSchema}] the columns actually read. Wider than the query
        needs means a \texttt{SELECT *} defeated column pruning.
  \item[\texttt{Batched: true}] vectorized columnar reads are enabled
        (\conf{spark.sql.parquet.enableVectorizedReader}); \texttt{false} falls
        back to row-at-a-time decoding.
\end{description}

\section{Exchange operators}
\label{sec:c-exchange}

Every \texttt{Exchange} is a shuffle and a stage boundary: serialize, write to
local disk, transfer over the network, re-read on the destination. Counting them
is the fastest estimate of a query's cost.

\begin{table}[H]
\centering\footnotesize
\renewcommand{\arraystretch}{1.2}
\begin{tabular}{@{}p{3cm}p{4.4cm}p{6cm}@{}}
\toprule
\textbf{Partitioning} & \textbf{In the plan} & \textbf{Created by} \\
\midrule
Hash & \texttt{hashpartitioning(col, n)} & joins, \texttt{GROUP BY}, \texttt{repartition(n, col)} \\
Range & \texttt{rangepartitioning(col ASC, n)} & global \texttt{ORDER BY}, \texttt{repartitionByRange} \\
Round-robin & \texttt{roundrobinpartitioning(n)} & \texttt{repartition(n)} with no column \\
Single & \texttt{SinglePartition} & global aggregate with no group key, \texttt{coalesce(1)} \\
\bottomrule
\end{tabular}
\end{table}

\begin{lstlisting}[language=text,numbers=none,caption={A shuffle exchange},captionpos=b]
Exchange hashpartitioning(product_id, 200), ENSURE_REQUIREMENTS, [plan_id=42]
\end{lstlisting}

\texttt{ENSURE\_REQUIREMENTS} means the planner's \texttt{EnsureRequirements} rule
inserted this exchange because the operator above it needs its input distributed
that way. \texttt{ReuseExchange} in a plan means two identical exchanges were
deduplicated --- one shuffle feeds both consumers.

\texttt{BroadcastExchange} is not a shuffle. It collects one small side to the
driver and ships it whole to every executor; it appears only below the build side
of a \texttt{BroadcastHashJoin} or \texttt{BroadcastNestedLoopJoin}.

\section{Sort operators}

\begin{lstlisting}[language=text,numbers=none]
Sort [product_id ASC NULLS FIRST], false, 0
\end{lstlisting}

The boolean is \texttt{true} for a global sort (a single ordered output across all
partitions, preceded by a \texttt{rangepartitioning} exchange) and \texttt{false}
for a local per-partition sort. A \texttt{SortMergeJoin} always shows two local
sorts, one per input, between the exchanges and the join. Sorts need a sort
buffer in execution memory and spill when it is exhausted (\Cref{ch:execmem}).

\section{Join operators}
\label{sec:c-joins}

\Cref{ch:joins} covers how the planner chooses; this is how to recognize the
choice in a plan.

\paragraph{BroadcastHashJoin.} \texttt{BroadcastHashJoin [\ldots], BuildRight}
(or \texttt{BuildLeft}) with a \texttt{BroadcastExchange} below the build side and
\emph{no} exchange on the probe (large) side. The cheapest join: no shuffle of
the big table.

\paragraph{SortMergeJoin.} \texttt{SortMergeJoin [\ldots], Inner} with two child
branches, each a \texttt{Sort} over an \texttt{Exchange hashpartitioning}. The
default for two large inputs, and the most expensive pattern --- two full
shuffles plus two sorts.

\paragraph{ShuffledHashJoin.} \texttt{ShuffledHashJoin [\ldots], BuildRight} with
two \texttt{Exchange} nodes but \emph{no} sorts: one side is shuffled and built
into a hash table in memory. Cheaper than sort-merge when the build side fits and
the sort is the bottleneck; enabled by
\conf{spark.sql.join.preferSortMergeJoin}\,=\,\texttt{false} or a
\texttt{SHUFFLE\_HASH} hint.

\paragraph{BroadcastNestedLoopJoin.} A non-equi join condition (\texttt{order\_total >
credit\_limit}). $O(N \times M)$. Fine when one side is tiny; a warning sign on two
large tables --- usually a missing or non-equi join key.

\begin{lstlisting}[language=text,numbers=none,caption={BroadcastNestedLoopJoin on a non-equi condition},captionpos=b]
BroadcastNestedLoopJoin BuildRight, Inner, (order_total#22 > credit_limit#41)
:- *(1) BatchScan orders
+- BroadcastExchange IdentityBroadcastMode
   +- *(2) Filter (segment#39 = VIP)
      +- *(2) BatchScan dim_customer
\end{lstlisting}

There is no \texttt{Exchange hashpartitioning} anywhere in this plan ---
the giveaway that the join key is not an equality. \texttt{dim\_customer}
itself is too large to broadcast by default (\Cref{sec:refenv-orders}), but
here it is filtered down to its small \texttt{VIP} segment first, and it is
\emph{that} filtered result, small enough to broadcast, that the
\texttt{BroadcastExchange} ships whole; every task on the \texttt{orders}
side then compares each of its rows against every row of that broadcast
copy, checking \texttt{order\_total > credit\_limit} directly. This is only
affordable because the broadcast side is small; the same plan shape with
neither side small enough to broadcast falls to \texttt{CartesianProduct}
instead.

\paragraph{CartesianProduct.} Every row of one side with every row of the other.
Unless it is a deliberate cross join, the \texttt{ON} clause is missing.

\section{Aggregation operators}
\label{sec:c-agg}

Aggregation is two-phase: a partial aggregate on each partition (map-side
combine), a shuffle, then a final aggregate.

\begin{lstlisting}[language=text,numbers=none,caption={Two-phase aggregation},captionpos=b]
*(3) HashAggregate(keys=[product_id], functions=[avg(on_time)])
+- Exchange hashpartitioning(product_id, 200), ENSURE_REQUIREMENTS
   +- *(2) HashAggregate(keys=[product_id], functions=[partial_avg(on_time)])
      +- ...
\end{lstlisting}

Reading bottom-up: the lower \texttt{partial\_avg(on\_time)} runs once per
input partition, producing one partial running-sum-and-count per
\texttt{product\_id} it saw locally; the \texttt{Exchange} then redistributes
those small partial results by \texttt{product\_id}, so every partial result
for the same product lands together; and the top \texttt{avg(on\_time)}
combines the partials that arrive for each key into the query's final
per-product average. Nothing about the specimen's real values is shuffled
directly --- only the already-shrunk partial aggregates are.

\begin{table}[H]
\centering\footnotesize
\renewcommand{\arraystretch}{1.2}
\begin{tabular}{@{}p{3.4cm}p{4.6cm}p{5.4cm}@{}}
\toprule
\textbf{Operator} & \textbf{Chosen when} & \textbf{Memory} \\
\midrule
\texttt{HashAggregate} & all aggregate buffers are fixed-width (numeric, date, timestamp) & off-heap hash map; falls back to sort-based under memory pressure \\
\texttt{ObjectHashAggregate} & a buffer holds a variable-width or complex type (\texttt{collect\_list}, \texttt{collect\_set}, some UDAFs) & on-heap objects; falls back to sort-based past \conf{spark.sql.objectHashAggregate.sortBased.fallbackThreshold} (128) distinct keys \\
\texttt{SortAggregate} & hash aggregation is not feasible & needs input pre-sorted by the group keys; handles any size, slower \\
\bottomrule
\end{tabular}
\end{table}

\section{Filter and Project}

\texttt{Filter} applies a predicate; \texttt{Project} selects columns and
evaluates expressions. Both are usually fused into a neighbouring scan or join
inside a \texttt{*(n)} block. A standalone \texttt{Filter} directly above a scan
is worth a look: if its predicate is not also in the scan's
\texttt{PushedFilters}, the pushdown failed and the predicate is being evaluated
row-by-row after reading (\Cref{sec:c-pushdown}).

\section{Whole-stage codegen markers}
\label{sec:c-codegen}

\begin{lstlisting}[language=text,numbers=none]
*(2) HashAggregate(keys=[city], functions=[sum(revenue)])
+- Exchange hashpartitioning(city, 200)
   +- *(1) HashAggregate(keys=[city], functions=[partial_sum(revenue)])
      +- *(1) Filter isnotnull(city)
         +- *(1) ColumnarToRow
            +- BatchScan sales [city, revenue]
\end{lstlisting}

\texttt{*(1)} is one generated Java method: scan, columnar-to-row, filter, and
partial aggregate run without leaving it. The \texttt{Exchange} breaks the chain
--- you cannot fuse across a network boundary --- so the final aggregate is
\texttt{*(2)}, a second generated method on the post-shuffle data. The number is
a sequential id, not an order or a priority. Operators with no star ran through
the interpreter; \Cref{sec:wholestage-codegen} lists why a stage loses codegen.

\section{AQE nodes in a final plan}
\label{sec:c-aqe}

When Adaptive Query Execution is on (\Cref{ch:aqe}), the plan the SQL tab shows
after the query finishes is rooted at \texttt{AdaptiveSparkPlan
isFinalPlan=true}. \texttt{isFinalPlan=false} is the compile-time plan from
\texttt{explain()} before an action --- it may change substantially at runtime.

\begin{description}
  \item[\texttt{AQEShuffleRead coalesced}] AQE merged small post-shuffle
        partitions after seeing their real sizes.
  \item[\texttt{AQEShuffleRead (skew=true)}] AQE split one or more oversized
        partitions.
  \item[\texttt{BroadcastQueryStage}] a \texttt{SortMergeJoin} was converted to a
        broadcast at runtime because the shuffled side turned out small.
  \item[\texttt{SortMergeJoin [\ldots], isSkew=true}] AQE detected and handled
        skew in this join.
\end{description}

Here is one final plan showing several of these together, for a join between
\texttt{orders} and the genuinely large \texttt{dim\_customer}
(\Cref{sec:refenv-orders}) during a week hit by the viral product's skew ---
both sides too large to broadcast, so this is a real sort-merge join, not one
AQE later converts:

\begin{lstlisting}[language=text,numbers=none,caption={A final AQE plan: a skew-split join feeding a coalesced aggregation},captionpos=b]
AdaptiveSparkPlan isFinalPlan=true
+- *(3) HashAggregate(keys=[category], functions=[avg(on_time)])
   +- AQEShuffleRead coalesced
      +- ShuffleQueryStage 2
         +- *(2) SortMergeJoin [product_id], [product_id], Inner, isSkew=true
            :- AQEShuffleRead (skew=true)
            :  +- ShuffleQueryStage 0
            +- AQEShuffleRead (skew=true)
               +- ShuffleQueryStage 1
\end{lstlisting}

The root confirms this is the plan as it actually ran, not the compile-time
guess. Both join inputs (\texttt{ShuffleQueryStage 0} and~\texttt{1}) carry
\texttt{AQEShuffleRead (skew=true)}: AQE found the viral product's oversized
partition after the shuffle and split it into several sub-partitions on the
\texttt{orders} side, and replicated the corresponding \texttt{dim\_customer}
partition to match, exactly \Cref{sec:aqe-skew} described. The join node's
own \texttt{isSkew=true} flag confirms this handling applied specifically to
this join, not just somewhere upstream. Above the join,
\texttt{AQEShuffleRead coalesced} shows a separate, independent AQE decision:
the aggregation's own shuffle output, originally many small partitions,
merged back down toward the 128~MB target before the final
\texttt{HashAggregate} runs over it --- one query, two unrelated
runtime-measured corrections, each visible as its own marker in the same
final plan.

A \texttt{BroadcastQueryStage} node appears in a different scenario: recall
from \Cref{sec:aqe-broadcast} that the specimen's \texttt{orders} $\bowtie$
\texttt{dim\_product} join is planned as sort-merge (Iceberg has no
column-level statistics to broadcast it up front), but \texttt{dim\_product}
turns out small once its own shuffle is actually measured. In \emph{that}
final plan, the join operator itself is no longer \texttt{SortMergeJoin} at
all --- AQE rewrites it to \texttt{BroadcastHashJoin}, with
\texttt{BroadcastQueryStage} standing in for the exchange-and-sort pair that
would have run on the now-broadcast side, and \texttt{orders} reads through
with no shuffle on that side either. A join cannot show both a
skew-split shuffle read \emph{and} a broadcast promotion on the same side:
promoting to broadcast removes that side's shuffle entirely, and skew-split
handling has nothing left to split once there is no post-shuffle partition
to measure.

\section{Predicate-pushdown killers}
\label{sec:c-pushdown}

Predicates that cannot be pushed to Parquet, ORC, or Iceberg, so they run
row-by-row after the read (\Cref{ch:pruning}):

\begin{itemize}
  \item a UDF: \texttt{WHERE my\_udf(col) = v}
  \item a function or complex expression on the column: \texttt{WHERE
        length(city) > 5}, \texttt{WHERE a + b > 100}
  \item a \texttt{CAST} on the column side: \texttt{WHERE CAST(s AS INT) > 100}
  \item a non-deterministic function: \texttt{WHERE rand() > 0.5}
  \item a leading-wildcard \texttt{LIKE}: \texttt{WHERE name LIKE '\%smith'}
        (a trailing wildcard, \texttt{'smith\%'}, does push)
  \item an \texttt{OR} with any non-pushable branch --- the whole \texttt{OR} fails
\end{itemize}

\paragraph{Verification.} Run \texttt{EXPLAIN FORMATTED}, find the scan, and check
that each partition-column filter is in \texttt{PartitionFilters} and each
data-column filter is in \texttt{PushedFilters}. If one is missing, it hit a
killer above.

Take \texttt{WHERE CAST(product\_id AS INT) > 1000} against the specimen's
\texttt{orders}, where \texttt{product\_id} is a \texttt{string}
(\Cref{sec:refenv-orders}). The scan node's detail block reads:

\begin{lstlisting}[language=text,numbers=none,caption={A CAST silently killing pushdown},captionpos=b]
(1) BatchScan orders
    PushedFilters:    []
    DataFilters:      [(cast(product_id as int) > 1000)]
    ReadSchema:       struct<product_id,order_ts,...>
\end{lstlisting}

\texttt{PushedFilters} is empty --- nothing made it to the storage layer at
all --- while the exact same condition sits in \texttt{DataFilters}, meaning
every row of every file is read and decoded before this predicate is
checked row-by-row in Spark itself. Compare that to the clean version of
the same intent, \texttt{WHERE product\_id > '1000'} (a plain string
comparison, no \texttt{CAST}): the scan's \texttt{PushedFilters} would show
\texttt{[IsNotNull(product\_id), GreaterThan(product\_id,1000)]} instead,
and \texttt{DataFilters} would be empty --- the same logical filter, but now
resolved by Iceberg's own statistics before a row is ever decoded. The
\texttt{CAST} is what moved it from one list to the other.
